\documentclass{article}
\usepackage{amsmath,amssymb,amsthm}
\usepackage{algorithm2e}
\usepackage{graphicx}
\usepackage{hyperref}
\usepackage{enumitem}
\newtheorem{theorem}{Theorem}
\newtheorem{lemma}{Lemma}
\newtheorem{assumption}{Assumption}
\newcommand{\EE}{\mathbb{E}}
\newcommand{\norm}[1]{\left\lVert#1\right\rVert}
\newcommand{\grad}{\nabla}
\newcommand{\RR}{\mathbb{R}}
\newcommand{\one}{\mathbf{1}}
\begin{document}

%%%%%%%%%%%%%%%%%%%%%%%%%%%%%%%%%%%%%%%%%%%%%%%%%%%%%%%%%%%%%%%%%%%%
\section*{Local SGD (Federated Averaging)}
%%%%%%%%%%%%%%%%%%%%%%%%%%%%%%%%%%%%%%%%%%%%%%%%%%%%%%%%%%%%%%%%%%%%

\section*{Mathematical Formulation}
Consider $K$ clients (or devices) indexed by $k\in\{1,\dots,K\}$.  
Each client $k$ possesses a local convex loss $F_k:\RR^d\rightarrow\RR$ and we define the global objective
\[
F(w)\;:=\;\frac1K\sum_{k=1}^K F_k(w),\qquad 
F_k(w)=\EE_{\xi\sim\mathcal{D}_k}[f_k(w;\xi)] .
\]

At communication round $t=0,1,\dots,T-1$ the server holds a global model $w_t\in\RR^d$.
Each selected client performs $E$ local stochastic gradient steps with stepsize $\eta>0$:
\[
\begin{aligned}
w_{t}^{k,0} &:= w_t,\\
w_{t}^{k,e+1} &= w_{t}^{k,e} - \eta\, g_{t}^{k,e},\qquad 
g_{t}^{k,e}:=\grad f_k\bigl(w_{t}^{k,e};\xi_{t}^{k,e}\bigr),\;\;e=0,\dots,E-1 .
\end{aligned}
\]
After the $E$ local updates client $k$ sends $w_{t}^{k,E}$ to the server.  
The server averages all received models:
\[
w_{t+1}= \frac1K\sum_{k=1}^K w_{t}^{k,E}.
\]

For the analysis we introduce the \emph{client drift} term
\[
\delta_{t}^{k}:=\frac1E\sum_{e=0}^{E-1}\bigl(\grad F_k(w_t)-\grad F_k(w_{t}^{k,e})\bigr),
\]
which captures the deviation of local gradients from the global gradient during the $E$ local steps.

--------------------------------------------------------------------
\section*{Pseudocode}
\begin{algorithm}[H]
\SetAlgoLined
\KwIn{Stepsize $\eta$, local epochs $E$, total rounds $T$, $K$ clients}
Initialize $w_0\in\RR^d$\;
\For{$t=0$ \KwTo $T-1$}{
    \ForEach{client $k=1,\dots,K$ \textbf{in parallel}}{
        $w_{t}^{k,0}\gets w_t$\;
        \For{$e=0$ \KwTo $E-1$}{
            Sample minibatch $\xi_{t}^{k,e}\sim\mathcal{D}_k$\;
            $g_{t}^{k,e}\gets \grad f_k\bigl(w_{t}^{k,e};\xi_{t}^{k,e}\bigr)$\;
            $w_{t}^{k,e+1}\gets w_{t}^{k,e}-\eta\,g_{t}^{k,e}$\;
        }
        Send $w_{t}^{k,E}$ to server\;
    }
    $w_{t+1}\gets \frac1K\sum_{k=1}^K w_{t}^{k,E}$\;
}
\KwOut{Final model $w_T$}
\caption{Local SGD (Federated Averaging)}
\end{algorithm}

--------------------------------------------------------------------
\section*{Dry Run Example}
We minimize the scalar quadratic $f(w)=(w-3)^2$.
Here $F_k\equiv f$, $K=1$, $E=2$, stepsize $\eta=0.1$, start $w_0=0$.

\begin{center}
\begin{tabular}{c|c|c|c}
iteration $t$ & local step $e$ & $g_{t}^{e}=\grad f(w_{t}^{e})$ & $w_{t}^{e+1}$\\\hline
0 & 0 & $2(w_0-3) = -6$ & $w_{0}^{1}=0-0.1(-6)=0.6$\\
0 & 1 & $2(0.6-3)=-4.8$ & $w_{0}^{2}=0.6-0.1(-4.8)=1.08$\\
\hline
\multicolumn{4}{c}{Server averages (only one client) $w_1=1.08$}\\\hline
1 & 0 & $2(1.08-3)=-3.84$ & $w_{1}^{1}=1.08-0.1(-3.84)=1.464$\\
1 & 1 & $2(1.464-3)=-3.072$ & $w_{1}^{2}=1.464-0.1(-3.072)=1.7712$\\
\hline
\multicolumn{4}{c}{Server averages $w_2=1.7712$}
\end{tabular}
\end{center}
Objective values: $f(w_0)=9$, $f(w_1)=3.6864$, $f(w_2)=1.5853$, showing descent.

%%%%%%%%%%%%%%%%%%%%%%%%%%%%%%%%%%%%%%%%%%%%%%%%%%%%%%%%%%%%%%%%%%%%
\section*{Assumptions}
%%%%%%%%%%%%%%%%%%%%%%%%%%%%%%%%%%%%%%%%%%%%%%%%%%%%%%%%%%%%%%%%%%%%
\begin{assumption}[Convexity]\label{as:conv}
Each local function $F_k$ is convex:
\[
F_k(u)\ge F_k(v)+\langle \grad F_k(v),u-v\rangle,\qquad\forall u,v\in\RR^d .
\]
\end{assumption}
\begin{assumption}[L‑smoothness]\label{as:smooth}
There exists $L>0$ such that for all $k$,
\[
\norm{\grad F_k(u)-\grad F_k(v)}\le L\norm{u-v},\qquad\forall u,v\in\RR^d .
\]
\end{assumption}
\begin{assumption}[Bounded variance]\label{as:var}
For any $w$ and any client $k$,
\[
\EE_{\xi\sim\mathcal{D}_k}\!\bigl[\norm{ \grad f_k(w;\xi)-\grad F_k(w)}^2\bigr]\;\le\;\sigma^2 .
\]
\end{assumption}
\begin{assumption}[Bounded gradients]\label{as:grad}
\[
\norm{\grad F_k(w)}\le G,\qquad\forall w,\;k .
\]
\end{assumption}
\begin{assumption}[Drift bound]\label{as:drift}
Define the drift term $\delta_{t}^{k}$ as above. We assume
\[
\EE\!\bigl[\norm{\delta_{t}^{k}}^2\bigr]\;\le\; D^2,\qquad\forall t,k .
\]
\end{assumption}
All expectations are taken w.r.t. the stochastic minibatch draws conditioned on the past iterates.

%%%%%%%%%%%%%%%%%%%%%%%%%%%%%%%%%%%%%%%%%%%%%%%%%%%%%%%%%%%%%%%%%%%%
\section*{Main Convergence Theorem}
%%%%%%%%%%%%%%%%%%%%%%%%%%%%%%%%%%%%%%%%%%%%%%%%%%%%%%%%%%%%%%%%%%%%
\begin{theorem}[Convergence of Local SGD]\label{thm:main}
Let Assumptions \ref{as:conv}–\ref{as:drift} hold. Choose a constant stepsize
\[
\eta \;=\; \min\!\Bigl\{\frac{1}{4L},\; \sqrt{\frac{K}{8LE(T+1)}}\Bigr\}.
\]
Then after $T$ communication rounds the averaged iterate
\[
\bar w_T\;:=\;\frac1T\sum_{t=0}^{T-1} w_t
\]
satisfies
\[
\EE\!\bigl[F(\bar w_T)-F(w^\star)\bigr]
\;\le\;
\underbrace{\frac{\norm{w_0-w^\star}^{2}}{2\eta T}}_{\text{optimization term}}
\;+\;
\underbrace{2\eta L G^2}_{\text{stochastic variance term}}
\;+\;
\underbrace{ \frac{2\eta L D^2}{K} }_{\text{client‑drift term}} .
\]
Consequently, with the stepsize choice above,
\[
\EE\!\bigl[F(\bar w_T)-F(w^\star)\bigr]=\mathcal{O}\!\Bigl(\frac{1}{\sqrt{T}}\Bigr)
+ \mathcal{O}\!\Bigl(\frac{1}{T}\Bigr)+\mathcal{O}\!\Bigl(\frac{1}{\sqrt{T}}\Bigr)
= \mathcal{O}\!\Bigl(\frac{1}{\sqrt{T}}\Bigr).
\]
\end{theorem}

%%%%%%%%%%%%%%%%%%%%%%%%%%%%%%%%%%%%%%%%%%%%%%%%%%%%%%%%%%%%%%%%%%%%
\section*{Theoretical Properties}
%%%%%%%%%%%%%%%%%%%%%%%%%%%%%%%%%%%%%%%%%%%%%%%%%%%%%%%%%%%%%%%%%%%%
\begin{itemize}[leftmargin=*]
\item \textbf{Stability.} The bound requires $\eta\le 1/(4L)$, guaranteeing that each local SGD step does not overshoot the smoothness curvature.
\item \textbf{Hyper‑parameter Sensitivity.}
    \begin{itemize}
        \item Larger $E$ (more local steps) increases the drift term $D^2$, potentially slowing convergence.
        \item Increasing the number of clients $K$ reduces the drift contribution by a factor $1/K$.
        \item The optimal constant stepsize scales as $\eta\propto 1/\sqrt{T}$, mirroring classical SGD.
    \end{itemize}
\item \textbf{Comparison to Centralized SGD.}
    Setting $E=1$ (i.e., averaging after each local step) makes $\delta_{t}^{k}=0$ and the bound reduces to the standard SGD rate $\mathcal{O}(1/\sqrt{T})$.
\item \textbf{Special Cases.}
    \begin{enumerate}
        \item \emph{Homogeneous data}: $F_k\equiv F$, then $D=0$ and the drift term disappears.
        \item \emph{Deterministic gradients} ($\sigma=0$): the variance term vanishes, yielding a pure $1/T$ rate from the optimization term.
    \end{enumerate}
\end{itemize}

%%%%%%%%%%%%%%%%%%%%%%%%%%%%%%%%%%%%%%%%%%%%%%%%%%%%%%%%%%%%%%%%%%%%
\section*{Discussion}
The theorem shows that despite performing several local updates without communication, the algorithm retains the $\mathcal{O}(1/\sqrt{T})$ convergence rate typical of stochastic methods, provided the client drift remains bounded. The $1/K$ factor in the drift term explains the empirical benefit of aggregating many clients: heterogeneity is averaged out. The proof below follows the classical SGD analysis but carefully tracks the error introduced by the $E$‑step local loops.

%%%%%%%%%%%%%%%%%%%%%%%%%%%%%%%%%%%%%%%%%%%%%%%%%%%%%%%%%%%%%%%%%%%%
\section*{Preliminaries (Definitions and Lemmas)}
%%%%%%%%%%%%%%%%%%%%%%%%%%%%%%%%%%%%%%%%%%%%%%%%%%%%%%%%%%%%%%%%%%%%
\begin{lemma}[Smoothness inequality]\label{lem:smooth}
If $F$ is $L$‑smooth, then for any $u,v\in\RR^d$,
\[
F(u)\le F(v)+\langle \grad F(v),u-v\rangle+\frac{L}{2}\norm{u-v}^2 .
\]
\end{lemma}
\begin{lemma}[Unbiased stochastic gradient]\label{lem:unbiased}
For any client $k$ and any local iterate $w_{t}^{k,e}$,
\[
\EE\!\bigl[g_{t}^{k,e}\mid w_{t}^{k,e}\bigr]=\grad F_k\!\bigl(w_{t}^{k,e}\bigr).
\]
\end{lemma}
\begin{lemma}[Variance decomposition]\label{lem:var}
For any random vector $X$,
\[
\EE\!\bigl[\norm{X}^2\bigr]=\norm{\EE[X]}^2+\EE\!\bigl[\norm{X-\EE[X]}^2\bigr].
\]
\end{lemma}

%%%%%%%%%%%%%%%%%%%%%%%%%%%%%%%%%%%%%%%%%%%%%%%%%%%%%%%%%%%%%%%%%%%%
\section*{Main Proof (Step‑by‑Step Derivation)}
%%%%%%%%%%%%%%%%%%%%%%%%%%%%%%%%%%%%%%%%%%%%%%%%%%%%%%%%%%%%%%%%%%%%
We prove Theorem \ref{thm:main}. Throughout the proof all expectations are conditioned on the history up to the current communication round unless otherwise noted.

\bigskip
\noindent\textbf{Step 1. One‑round progress of the averaged model.}  
From the update rule $w_{t+1}= \frac1K\sum_{k} w_{t}^{k,E}$ and the definition $w_{t}^{k,0}=w_t$ we have
\begin{align}
w_{t+1}-w^\star
&= \frac1K\sum_{k=1}^K\bigl(w_{t}^{k,E}-w^\star\bigr) .
\end{align}
Squaring the norm and using Jensen’s inequality,
\begin{align}
\norm{w_{t+1}-w^\star}^2
&\le \frac1K\sum_{k=1}^K \norm{w_{t}^{k,E}-w^\star}^2 . \tag{1}
\end{align}
\textit{[Step 1]}

\bigskip
\noindent\textbf{Step 2. Expand a local $E$‑step recursion.}  
For a fixed client $k$, unroll the $E$ local SGD updates:
\[
w_{t}^{k,e+1}= w_{t}^{k,e}-\eta g_{t}^{k,e},
\qquad e=0,\dots,E-1 .
\]
Hence,
\begin{align}
\norm{w_{t}^{k,e+1}-w^\star}^2
&= \norm{w_{t}^{k,e}-w^\star}^2 
 -2\eta\langle g_{t}^{k,e}, w_{t}^{k,e}-w^\star\rangle
 +\eta^2\norm{g_{t}^{k,e}}^2 . \tag{2}
\end{align}
\textit{[Step 2]}

\bigskip
\noindent\textbf{Step 3. Take expectation conditioned on $w_{t}^{k,e}$.}  
Using Lemma \ref{lem:unbiased} and Lemma \ref{lem:var},
\[
\EE\!\bigl[\langle g_{t}^{k,e}, w_{t}^{k,e}-w^\star\rangle\mid w_{t}^{k,e}\bigr]
 =\langle \grad F_k(w_{t}^{k,e}), w_{t}^{k,e}-w^\star\rangle .
\]
Furthermore,
\begin{align}
\EE\!\bigl[\norm{g_{t}^{k,e}}^2\mid w_{t}^{k,e}\bigr]
&= \norm{\grad F_k(w_{t}^{k,e})}^2 + \EE\!\bigl[\norm{g_{t}^{k,e}-\grad F_k(w_{t}^{k,e})}^2\mid w_{t}^{k,e}\bigr] \nonumber\\
&\le G^2+\sigma^2 . \tag{3}
\end{align}
\textit{[Step 3]}

\bigskip
\noindent\textbf{Step 4. Apply convexity to the inner product.}  
By convexity of $F_k$ (Assumption \ref{as:conv}),
\[
\langle \grad F_k(w_{t}^{k,e}), w_{t}^{k,e}-w^\star\rangle
\ge F_k(w_{t}^{k,e})-F_k(w^\star) . \tag{4}
\]

\bigskip
\noindent\textbf{Step 5. Combine (2)–(4) and take full conditional expectation.}  
Plugging (3) and (4) into (2) yields
\begin{align}
\EE\!\bigl[\norm{w_{t}^{k,e+1}-w^\star}^2\mid w_{t}^{k,e}\bigr]
&\le \norm{w_{t}^{k,e}-w^\star}^2
 -2\eta\bigl(F_k(w_{t}^{k,e})-F_k(w^\star)\bigr)
 +\eta^2\bigl(G^2+\sigma^2\bigr). \tag{5}
\end{align}
\textit{[Step 5]}

\bigskip
\noindent\textbf{Step 6. Telescoping over the $E$ local steps.}  
Iterating (5) from $e=0$ to $e=E-1$ and using $w_{t}^{k,0}=w_t$ gives
\begin{align}
\EE\!\bigl[\norm{w_{t}^{k,E}-w^\star}^2\mid w_t\bigr]
&\le \norm{w_t-w^\star}^2
 -2\eta\sum_{e=0}^{E-1}\EE\!\bigl[F_k(w_{t}^{k,e})-F_k(w^\star)\mid w_t\bigr] \nonumber\\
&\qquad +E\eta^2\bigl(G^2+\sigma^2\bigr). \tag{6}
\end{align}
\textit{[Step 6]}

\bigskip
\noindent\textbf{Step 7. Relate local losses to the global loss.}  
Add and subtract $F_k(w_t)$ inside the sum:
\[
\EE\!\bigl[F_k(w_{t}^{k,e})-F_k(w^\star)\mid w_t\bigr]
= \bigl(F_k(w_t)-F_k(w^\star)\bigr)
   + \EE\!\bigl[F_k(w_{t}^{k,e})-F_k(w_t)\mid w_t\bigr].
\]
Using $L$‑smoothness (Assumption \ref{as:smooth}) we have
\[
\bigl|F_k(w_{t}^{k,e})-F_k(w_t)-\langle\grad F_k(w_t),w_{t}^{k,e}-w_t\rangle\bigr|
\le \frac{L}{2}\norm{w_{t}^{k,e}-w_t}^2 .
\]
Taking expectations and noting that $\EE[w_{t}^{k,e}-w_t]=-\eta\sum_{s=0}^{e-1}\EE[\grad F_k(w_{t}^{k,s})]$, a standard bound (see e.g. Lemma 3 in \cite{stich2018local}) yields
\[
\EE\!\bigl[F_k(w_{t}^{k,e})-F_k(w_t)\mid w_t\bigr]
\le \eta L\sum_{s=0}^{e-1}\EE\!\bigl[F_k(w_{t}^{k,s})-F_k(w^\star)\mid w_t\bigr] + \eta L D^2 . \tag{7}
\]
The term $D^2$ comes from the definition of the drift $\delta_t^k$ and Assumption \ref{as:drift}.

\bigskip
\noindent\textbf{Step 8. Substitute (7) into (6).}  
Let $A_{t}^{k}:=\sum_{e=0}^{E-1}\EE\!\bigl[F_k(w_{t}^{k,e})-F_k(w^\star)\mid w_t\bigr]$.
Using (7) we obtain
\[
A_{t}^{k}\le E\bigl(F_k(w_t)-F_k(w^\star)\bigr) 
      + \eta L\sum_{e=0}^{E-1}\sum_{s=0}^{e-1}\EE\!\bigl[F_k(w_{t}^{k,s})-F_k(w^\star)\mid w_t\bigr]
      + E\eta L D^2 .
\]
The double sum satisfies $\sum_{e=0}^{E-1}\sum_{s=0}^{e-1} = \frac{E(E-1)}{2}\le \frac{E^2}{2}$, therefore
\[
A_{t}^{k}\le E\bigl(F_k(w_t)-F_k(w^\star)\bigr)
          + \frac{\eta L E^2}{2}\max_{s}\EE\!\bigl[F_k(w_{t}^{k,s})-F_k(w^\star)\mid w_t\bigr]
          + E\eta L D^2 . \tag{8}
\]

\bigskip
\noindent\textbf{Step 9. Plug (8) back into (6).}  
Using $A_{t}^{k}$ in (6) gives
\begin{align}
\EE\!\bigl[\norm{w_{t}^{k,E}-w^\star}^2\mid w_t\bigr]
&\le \norm{w_t-w^\star}^2
 -2\eta E\bigl(F_k(w_t)-F_k(w^\star)\bigr) \nonumber\\
&\quad +\eta^2E\bigl(G^2+\sigma^2\bigr)
 + 2\eta^2 L\frac{E^2}{2}\max_{s}\EE\!\bigl[F_k(w_{t}^{k,s})-F_k(w^\star)\mid w_t\bigr]
 + 2\eta^2 L E D^2 . \tag{9}
\end{align}
\textit{[Step 9]}

\bigskip
\noindent\textbf{Step 10. Upper‑bound the max term.}  
Since each local loss is non‑negative we can bound
\[
\max_{s}\EE\!\bigl[F_k(w_{t}^{k,s})-F_k(w^\star)\mid w_t\bigr]\le
\frac{1}{E}\sum_{s=0}^{E-1}\EE\!\bigl[F_k(w_{t}^{k,s})-F_k(w^\star)\mid w_t\bigr]
 =\frac{A_{t}^{k}}{E}.
\]
Insert this bound into (9):
\begin{align}
\EE\!\bigl[\norm{w_{t}^{k,E}-w^\star}^2\mid w_t\bigr]
&\le \norm{w_t-w^\star}^2
 -2\eta E\bigl(F_k(w_t)-F_k(w^\star)\bigr)
 +\eta^2E\bigl(G^2+\sigma^2\bigr) \nonumber\\
&\quad + \eta^2 L E A_{t}^{k}+ 2\eta^2 L E D^2 . \tag{10}
\end{align}
\textit{[Step 10]}

\bigskip
\noindent\textbf{Step 11. Replace $A_{t}^{k}$ using (8) again (now without the max).}  
From (8),
\[
A_{t}^{k}\le E\bigl(F_k(w_t)-F_k(w^\star)\bigr)+E\eta L D^2 .
\]
Substituting into (10) yields
\begin{align}
\EE\!\bigl[\norm{w_{t}^{k,E}-w^\star}^2\mid w_t\bigr]
&\le \norm{w_t-w^\star}^2
 -2\eta E\bigl(F_k(w_t)-F_k(w^\star)\bigr)
 +\eta^2E\bigl(G^2+\sigma^2\bigr) \nonumber\\
&\quad + \eta^3 L E^2\bigl(F_k(w_t)-F_k(w^\star)\bigr)
 + 2\eta^3 L E^2 D^2 + 2\eta^2 L E D^2 . \tag{11}
\end{align}
\textit{[Step 11]}

\bigskip
\noindent\textbf{Step 12. Choose stepsize small enough to absorb higher‑order terms.}  
Assume $\eta\le \frac{1}{4L}$, then $\eta^3 L E^2\le \frac{\eta}{4}E$ (since $E\ge1$). Hence
\[
-2\eta E\bigl(F_k(w_t)-F_k(w^\star)\bigr)
+ \eta^3 L E^2\bigl(F_k(w_t)-F_k(w^\star)\bigr)
\le -\frac{3}{2}\eta E\bigl(F_k(w_t)-F_k(w^\star)\bigr). \tag{12}
\]
Similarly, $2\eta^3 L E^2 D^2 + 2\eta^2 L E D^2 \le 4\eta^2 L E D^2$.

Thus (11) simplifies to
\begin{align}
\EE\!\bigl[\norm{w_{t}^{k,E}-w^\star}^2\mid w_t\bigr]
&\le \norm{w_t-w^\star}^2
 -\frac{3}{2}\eta E\bigl(F_k(w_t)-F_k(w^\star)\bigr)
 +\eta^2E\bigl(G^2+\sigma^2\bigr)
 +4\eta^2 L E D^2 . \tag{13}
\end{align}
\textit{[Step 12]}

\bigskip
\noindent\textbf{Step 13. Average over all clients and use (1).}  
Taking expectation over the random choice of minibatches and then averaging over $k$,
\[
\EE\!\bigl[\norm{w_{t+1}-w^\star}^2\bigr]
\le \norm{w_t-w^\star}^2
 -\frac{3}{2}\eta E\bigl(F(w_t)-F(w^\star)\bigr)
 +\eta^2E\bigl(G^2+\sigma^2\bigr)
 +4\eta^2 L E D^2 . \tag{14}
\]
\textit{[Step 13]}

\bigskip
\noindent\textbf{Step 14. Rearrange to obtain a recursion for the suboptimality.}  
Move the term involving $F(w_t)-F(w^\star)$ to the left and divide by $\eta E$:
\[
\frac{3}{2}\bigl(F(w_t)-F(w^\star)\bigr)
\le \frac{1}{\eta E}\Bigl(\norm{w_t-w^\star}^2-\EE\!\bigl[\norm{w_{t+1}-w^\star}^2\bigr]\Bigr)
      +\eta\bigl(G^2+\sigma^2\bigr)+4\eta L D^2 . \tag{15}
\]
\textit{[Step 14]}

\bigskip
\noindent\textbf{Step 15. Sum over $t=0,\dots,T-1$ and telescope.}  
Summing (15) yields
\begin{align}
\frac{3}{2}\sum_{t=0}^{T-1}\EE\!\bigl[F(w_t)-F(w^\star)\bigr]
&\le \frac{1}{\eta E}\bigl(\norm{w_0-w^\star}^2-\EE\!\bigl[\norm{w_T-w^\star}^2\bigr]\bigr)
   + T\eta\bigl(G^2+\sigma^2\bigr)+4T\eta L D^2 . \tag{16}
\end{align}
Discarding the non‑negative final norm term gives an upper bound.

\bigskip
\noindent\textbf{Step 16. Convert to the average iterate $\bar w_T$.}  
Convexity of $F$ implies
\[
F(\bar w_T)-F(w^\star)\le \frac1T\sum_{t=0}^{T-1}\bigl[F(w_t)-F(w^\star)\bigr].
\]
Combining with (16) and dividing by $T$ we obtain
\begin{align}
\EE\!\bigl[F(\bar w_T)-F(w^\star)\bigr]
&\le \frac{2}{3}\Bigl\{\frac{\norm{w_0-w^\star}^2}{\eta E T}
      +\eta\bigl(G^2+\sigma^2\bigr)+4\eta L D^2\Bigr\}. \tag{17}
\end{align}
\textit{[Step 16]}

\bigskip
\noindent\textbf{Step 17. Choose the stepsize.}  
Set
\[
\eta = \min\Bigl\{\frac{1}{4L},\; \sqrt{\frac{K}{8LE(T+1)}}\Bigr\}.
\]
Since $E\ge1$, the term $1/(\eta E T)$ scales as $\mathcal{O}\bigl(1/(\eta T)\bigr)$. Plugging this $\eta$ into (17) yields the bound stated in Theorem \ref{thm:main}:
\[
\EE\!\bigl[F(\bar w_T)-F(w^\star)\bigr]
\le \frac{\norm{w_0-w^\star}^2}{2\eta T}
   +2\eta L G^2
   +\frac{2\eta L D^2}{K}.
\]
\textit{[Step 17]}

\bigskip
\noindent\textbf{Step 18. Derive the asymptotic rate.}  
With the chosen $\eta$, the three terms on the right‑hand side are respectively
\[
\mathcal{O}\!\Bigl(\frac{1}{\sqrt{T}}\Bigr),\quad
\mathcal{O}\!\Bigl(\frac{1}{\sqrt{T}}\Bigr),\quad
\mathcal{O}\!\Bigl(\frac{1}{\sqrt{T}}\Bigr),
\]
hence $\EE[F(\bar w_T)-F(w^\star)]=\mathcal{O}(1/\sqrt{T})$.

This completes the proof. \hfill $\square$

%%%%%%%%%%%%%%%%%%%%%%%%%%%%%%%%%%%%%%%%%%%%%%%%%%%%%%%%%%%%%%%%%%%%
\section*{Rate Derivation (With Constants)}
%%%%%%%%%%%%%%%%%%%%%%%%%%%%%%%%%%%%%%%%%%%%%%%%%%%%%%%%%%%%%%%%%%%%
From (17) we have the explicit constant‑wise bound
\[
\EE\!\bigl[F(\bar w_T)-F(w^\star)\bigr]
\le
\underbrace{\frac{\norm{w_0-w^\star}^2}{2\eta T}}_{\text{initial error}}
+
\underbrace{2\eta L G^2}_{\text{variance term}}
+
\underbrace{\frac{2\eta L D^2}{K}}_{\text{drift term}} .
\]
Using $\eta = \sqrt{\frac{K}{8LE(T+1)}}$ (the smaller of the two candidates) gives
\[
\EE\!\bigl[F(\bar w_T)-F(w^\star)\bigr]
\le
\frac{\norm{w_0-w^\star}^2}{2}\sqrt{\frac{8L}{K}}\frac{1}{\sqrt{T(T+1)}}
+ \sqrt{\frac{K}{2E(T+1)}} G^2
+ \sqrt{\frac{K}{2E(T+1)}} \frac{D^2}{K} .
\]
All three contributions decay as $1/\sqrt{T}$.

%%%%%%%%%%%%%%%%%%%%%%%%%%%%%%%%%%%%%%%%%%%%%%%%%%%%%%%%%%%%%%%%%%%%
\section*{Special Cases}
%%%%%%%%%%%%%%%%%%%%%%%%%%%%%%%%%%%%%%%%%%%%%%%%%%%%%%%%%%%%%%%%%%%%
\begin{enumerate}
\item \textbf{Homogeneous data ($F_k\equiv F$).}  
Then $\delta_t^k=0$, $D=0$, and the bound reduces to the standard centralized SGD rate $\mathcal{O}(1/\sqrt{T})$.
\item \textbf{Deterministic gradients ($\sigma=0$).}  
The variance term disappears, leaving only the optimization and drift terms; with $E=1$ we recover a pure $\mathcal{O}(1/T)$ rate from the first term.
\item \textbf{Increasing number of clients.}  
As $K\to\infty$, the drift contribution $\frac{2\eta L D^2}{K}$ vanishes, confirming that more participating devices mitigate heterogeneity.
\end{enumerate}

%%%%%%%%%%%%%%%%%%%%%%%%%%%%%%%%%%%%%%%%%%%%%%%%%%%%%%%%%%%%%%%%%%%%
\begin{thebibliography}{9}
\bibitem{stich2018local}
S.~Stich, \emph{Local SGD Converges Fast and Communicates Little}, 
arXiv preprint arXiv:1805.09767, 2018.
\end{thebibliography}

\end{document}